\begin{lemma}
Let $n\ge1$ be an integer, let $1<p<\infty$, let $k\ge n$ be an integer and let
$A\colon\mathbb{Z}\to\mathbb{C}$.  Write $\lambda_k=\lambda_k(n,p)$ for the
exponent of \noderef{LambdaCoef}, $\epsilon_k=\epsilon_k(n,p)$ for the
coefficient of \noderef{EpsCoef}, and $T_k$ for the $k$-th summand of (19) in
the sense of \noderef{Thm2Term}, so that
\[
  T_k(\xi)=\bigl(A_k\xi^k+\epsilon_k\overline{A_k}\bar\xi^{\,k}\bigr)
    |\xi|^{\lambda_k+n-k}\xi^{-n}.
\]
Then:
\begin{enumerate}
\item[(a)] $T_k$ is $C^\infty$ at every $\xi\neq0$;
\item[(b)] for every $r>0$ and every $\theta\in\mathbb{R}$,
\[
  T_k\bigl(re^{i\theta}\bigr)
  =r^{\lambda_k}\Bigl(A_ke^{i(k-n)\theta}
    +\epsilon_k\overline{A_k}e^{i(-k-n)\theta}\Bigr),
\]
so $T_k$ is homogeneous of degree $\lambda_k$ with angular profile a
trigonometric polynomial of frequencies $k-n$ and $-k-n$;
\item[(c)] $T_k$ solves equation (18) on the unit disk
$\mathbb{B}=\{\xi:|\xi|<1\}$, in the sense of \noderef{SolvesEq18}.
\end{enumerate}

Part (c) is the paper's assertion at line 357 that ``our arguments can be
turned back'', and it is the whole content of the converse half of Theorem 2.
No restriction on $p$ is imposed; in particular (c) holds at $p=2$, where
$\lambda_k=k-n$, $\epsilon_k=0$ and $T_k(\xi)=A_k\xi^{k-n}$ is holomorphic.
\end{lemma}
\begin{proof}
Throughout, write $L=\lambda_k$, $e=\epsilon_k$, $\mu=L+n-k$ and
\[
  F(\xi)=A_k\xi^k+e\,\overline{A_k}\,\bar\xi^{\,k},
\]
so that by \noderef{Thm2Term}
\begin{equation}\tag{0}
  T_k(\xi)=F(\xi)\,|\xi|^{\mu}\,\xi^{-n}\qquad(\xi\neq0).
\end{equation}
Note $e\in\mathbb{R}$ by \noderef{EpsCoef}.

\emph{Step 0: $L\ge0$, $L+n+k>0$, and the indicial equation.}  By part (23) of
\noderef{LambdaEpsEstimates}, $\lambda_n(n,p)=0$, and by part (25) of the same
node $\lambda_k\ge(p-1)/(np)>0$ for $k\ge n+1$.  Since $k\ge n$, in either case
\begin{equation}\tag{1}
  L\ge0,\qquad\text{hence}\qquad L+n+k\ge0+1+1>0 ,
\end{equation}
using $n\ge1$ and $k\ge n\ge1$.

Next the indicial equation.  By \noderef{LambdaCoef}, with $|k|=k$,
\[
  2L+pn=\sqrt{4k^2(p-1)+n^2(p-2)^2},
\]
the right-hand side being a nonnegative real number.  Squaring,
$(2L+pn)^2=4k^2(p-1)+n^2(p-2)^2$, that is
$4L^2+4pnL+n^2p^2=4k^2(p-1)+n^2(p-2)^2$; hence, using
$(p-2)^2-p^2=-4(p-1)$,
\[
  4L^2+4pnL=4k^2(p-1)+n^2\bigl[(p-2)^2-p^2\bigr]=4(k^2-n^2)(p-1),
\]
and dividing by $4$,
\begin{equation}\tag{I}
  L^2+pnL+(n^2-k^2)(p-1)=0 .
\end{equation}
This is the indicial equation at paper.tex:315.  It is the only property of $L$
used below, so formula (37)'s square root is never unfolded again.

\emph{Step 1: proof of (b).}  Let $r>0$, $\theta\in\mathbb{R}$ and
$\xi=re^{i\theta}$.  Since $|e^{i\theta}|=1$ we have $|\xi|=r>0$, and
\[
  \xi^k=r^ke^{ik\theta},\qquad \bar\xi^{\,k}=r^ke^{-ik\theta},\qquad
  \xi^{-n}=r^{-n}e^{-in\theta},
\]
the last because $\xi\neq0$.  Substituting into (0) and using $r^ar^b=r^{a+b}$
for $r>0$,
\[
  T_k(\xi)=\bigl(A_kr^ke^{ik\theta}+e\overline{A_k}r^ke^{-ik\theta}\bigr)
    r^{\mu}r^{-n}e^{-in\theta}
  =r^{k+\mu-n}\bigl(A_ke^{i(k-n)\theta}+e\overline{A_k}e^{i(-k-n)\theta}\bigr).
\]
The exponent is $k+\mu-n=k+(L+n-k)-n=L$, which is (b).

\emph{Step 2: proof of (a).}  Fix $\xi\neq0$.  The four factors in (0) are
$C^\infty$ at $\xi$:
\begin{itemize}
\item $z\mapsto z^{k}$ is a polynomial, hence $C^\infty$ on $\mathbb{C}$;
\item $z\mapsto\bar z$ is $\mathbb{R}$-linear, hence $C^\infty$, so
  $z\mapsto\bar z^{\,k}$ is $C^\infty$ on $\mathbb{C}$;
\item $z\mapsto|z|^{\mu}$ is $C^\infty$ at $\xi$ by part (a) of
  \noderef{WirtingerModulusPowerRule};
\item $z\mapsto z^{-n}$ is $C^\infty$ on $\mathbb{C}-\{0\}$, being the
  reciprocal of the nonvanishing polynomial $z^{n}$ there.
\end{itemize}
A product of finitely many functions $C^\infty$ at $\xi$ is $C^\infty$ at
$\xi$, so $T_k$ is $C^\infty$ at $\xi$.  This is (a); in particular all the
first derivatives used in Steps 3--5 exist, and each factor is in particular
real-differentiable at $\xi$, which is the hypothesis the Wirtinger rules
require.

\emph{Step 3: the Wirtinger derivatives of the four factors.}  Let
$\Omega=\mathbb{C}-\{0\}$, which is open, and fix $\xi\in\Omega$.

The identity map $z\mapsto z$ has real Fréchet derivative the identity, so by
\noderef{DZ} and \noderef{DZbar}
\[
  \frac{\partial z}{\partial\xi}=\tfrac12\bigl(1-i\cdot i\bigr)=1,\qquad
  \frac{\partial z}{\partial\bar\xi}=\tfrac12\bigl(1+i\cdot i\bigr)=0 .
\]
The identity map is $C^\infty$ and nonvanishing on $\Omega$, so
\noderef{WirtingerPowerRule} applies to it with the integer exponents $k$ and
$-n$ and gives
\begin{equation}\tag{2}
  \frac{\partial\,\xi^{k}}{\partial\xi}=k\,\xi^{k-1},\qquad
  \frac{\partial\,\xi^{k}}{\partial\bar\xi}=0,\qquad
  \frac{\partial\,\xi^{-n}}{\partial\xi}=-n\,\xi^{-n-1},\qquad
  \frac{\partial\,\xi^{-n}}{\partial\bar\xi}=0 .
\end{equation}
By \noderef{WirtingerConjRule} applied to the identity map, the map
$c(z)=\bar z$ satisfies
$\partial c/\partial\xi=\overline{\partial z/\partial\bar\xi}=0$ and
$\partial c/\partial\bar\xi=\overline{\partial z/\partial\xi}=1$; and $c$ is
$C^\infty$ and nonvanishing on $\Omega$, so \noderef{WirtingerPowerRule}
applied to $c$ with exponent $k$ gives
\begin{equation}\tag{3}
  \frac{\partial\,\bar\xi^{\,k}}{\partial\xi}=k\,\bar\xi^{\,k-1}\cdot0=0,\qquad
  \frac{\partial\,\bar\xi^{\,k}}{\partial\bar\xi}=k\,\bar\xi^{\,k-1} .
\end{equation}
By \noderef{WirtingerModulusPowerRule},
\begin{equation}\tag{4}
  \frac{\partial\,|\xi|^{\mu}}{\partial\xi}=\frac{\mu}{2}\frac{|\xi|^{\mu}}{\xi},
  \qquad
  \frac{\partial\,|\xi|^{\mu}}{\partial\bar\xi}
    =\frac{\mu}{2}\frac{|\xi|^{\mu}}{\bar\xi}.
\end{equation}
From (2) and (3), by $\mathbb{C}$-linearity of each Wirtinger operator in the
function (immediate from \noderef{DZ}, \noderef{DZbar} and linearity of the
Fréchet derivative),
\begin{equation}\tag{5}
  \frac{\partial F}{\partial\xi}(\xi)=k\,A_k\xi^{k-1},\qquad
  \frac{\partial F}{\partial\bar\xi}(\xi)=k\,e\,\overline{A_k}\,\bar\xi^{\,k-1}.
\end{equation}

\emph{Step 4: the Wirtinger derivatives of $T_k$.}  Apply
\noderef{WirtingerProductRule} twice to the product (0) --- legitimate by
Step 2, since each factor is real-differentiable at $\xi$ --- and substitute
(2), (4), (5):
\begin{align*}
  \frac{\partial T_k}{\partial\xi}(\xi)
  &=\frac{\partial F}{\partial\xi}|\xi|^{\mu}\xi^{-n}
    +F\cdot\frac{\mu}{2}\frac{|\xi|^{\mu}}{\xi}\cdot\xi^{-n}
    +F|\xi|^{\mu}\cdot\bigl(-n\xi^{-n-1}\bigr)\\
  &=|\xi|^{\mu}\xi^{-n}\Bigl[k\,A_k\xi^{k-1}
      +\Bigl(\frac{\mu}{2}-n\Bigr)\frac{F(\xi)}{\xi}\Bigr],
\end{align*}
\begin{align*}
  \frac{\partial T_k}{\partial\bar\xi}(\xi)
  &=\frac{\partial F}{\partial\bar\xi}|\xi|^{\mu}\xi^{-n}
    +F\cdot\frac{\mu}{2}\frac{|\xi|^{\mu}}{\bar\xi}\cdot\xi^{-n}+0
  =|\xi|^{\mu}\xi^{-n}\Bigl[k\,e\,\overline{A_k}\,\bar\xi^{\,k-1}
      +\frac{\mu}{2}\frac{F(\xi)}{\bar\xi}\Bigr].
\end{align*}
Both displays used $\xi^{-n}/\xi=\xi^{-n-1}$ and $\xi\neq0$.
These two closed forms, with the factor $|\xi|^{\mu}\xi^{-n}/\xi$ and
$|\xi|^{\mu}\xi^{-n}/\bar\xi$ pulled out, are
\noderef{Thm2TermWirtingerDerivs}, which proves them with no constraint on
$n$, $p$ or $k$.

\emph{Step 5: reduction of (18) to two scalar identities.}  Fix $\xi\neq0$ and
put $c=\tfrac12-\tfrac1p$.  Equation (18) at $\xi$ reads
\begin{equation}\tag{18}
  \frac{\partial T_k}{\partial\bar\xi}
  =c\Bigl[\frac{\xi}{\bar\xi}\frac{\partial T_k}{\partial\xi}
    +\frac{\bar\xi^{\,n}}{\xi^{n}}
      \overline{\Bigl(\frac{\partial T_k}{\partial\xi}\Bigr)}\Bigr].
\end{equation}
Conjugating the first display of Step 4, and using that $|\xi|^{\mu}$ is real
and $e$ is real,
\[
  \overline{\frac{\partial T_k}{\partial\xi}(\xi)}
  =|\xi|^{\mu}\bar\xi^{\,-n}\Bigl[k\,\overline{A_k}\,\bar\xi^{\,k-1}
    +\Bigl(\frac{\mu}{2}-n\Bigr)\frac{\overline{F(\xi)}}{\bar\xi}\Bigr],
  \qquad
  \overline{F(\xi)}=\overline{A_k}\bar\xi^{\,k}+e\,A_k\xi^{k},
\]
so, since $\bar\xi^{\,n}\bar\xi^{\,-n}=1$,
\[
  \frac{\bar\xi^{\,n}}{\xi^{n}}\overline{\frac{\partial T_k}{\partial\xi}(\xi)}
  =|\xi|^{\mu}\xi^{-n}\Bigl[k\,\overline{A_k}\,\bar\xi^{\,k-1}
    +\Bigl(\frac{\mu}{2}-n\Bigr)\frac{\overline{F(\xi)}}{\bar\xi}\Bigr].
\]
The common factor $|\xi|^{\mu}\xi^{-n}$ is nonzero, since $|\xi|^{\mu}>0$ and
$\xi\neq0$.  Dividing (18) by it and then multiplying by $\bar\xi\neq0$, and
writing
\[
  u=A_k\xi^{k},\qquad v=\overline{A_k}\,\bar\xi^{\,k},
\]
so that $F=u+e\,v$, $\overline{F}=v+e\,u$ and $F+\overline{F}=(1+e)(u+v)$,
equation (18) becomes
\[
  k\,e\,v+\frac{\mu}{2}\bigl(u+e\,v\bigr)
  =c\Bigl[k\,u+k\,v+\Bigl(\frac{\mu}{2}-n\Bigr)(1+e)(u+v)\Bigr],
\]
where on the left $\bar\xi\cdot\overline{A_k}\bar\xi^{\,k-1}=v$ and
$\bar\xi\cdot F/\bar\xi=F$, and on the right
$\bar\xi\cdot(\xi/\bar\xi)\,A_k k\xi^{k-1}=k\,u$ and
$\bar\xi\cdot(\xi/\bar\xi)(\mu/2-n)F/\xi=(\mu/2-n)F$, while
$\bar\xi\cdot k\overline{A_k}\bar\xi^{\,k-1}=k\,v$ and
$\bar\xi\cdot(\mu/2-n)\overline{F}/\bar\xi=(\mu/2-n)\overline{F}$.  Grouping
the coefficients of $u$ and of $v$, (18) holds at $\xi$ as soon as the two
scalar identities
\begin{equation}\tag{$\alpha$}
  \frac{\mu}{2}=c\Bigl[k+\Bigl(\frac{\mu}{2}-n\Bigr)(1+e)\Bigr]
\end{equation}
and
\begin{equation}\tag{$\beta$}
  e\Bigl(k+\frac{\mu}{2}\Bigr)=c\Bigl[k+\Bigl(\frac{\mu}{2}-n\Bigr)(1+e)\Bigr]
\end{equation}
hold, because then the left side is $\tfrac{\mu}{2}u+e(k+\tfrac{\mu}{2})v$ and
the right side is the same bracket times $u+v$.  Note that no property of
$\xi$, $A_k$ or $n$ beyond $\xi\neq0$ entered this reduction: the remaining
work is pure algebra in $L$, $n$, $k$ and $p$.

\emph{Step 6: proof of $(\beta)$ from $(\alpha)$.}  By \noderef{EpsCoef} and
(1),
\[
  e=\frac{L+n-k}{L+n+k}=\frac{\mu}{L+n+k},
\]
the denominator being positive by (1).  Also
$k+\frac{\mu}{2}=k+\frac{L+n-k}{2}=\frac{L+n+k}{2}$.  Hence
\begin{equation}\tag{6}
  e\Bigl(k+\frac{\mu}{2}\Bigr)
  =\frac{\mu}{L+n+k}\cdot\frac{L+n+k}{2}=\frac{\mu}{2},
\end{equation}
so $(\beta)$ is exactly $(\alpha)$.  (This identity is an algebraic triviality:
it holds for every value of $L$, and uses neither (I) nor the sign of $p-2$.)

\emph{Step 7: proof of $(\alpha)$.}  First simplify the bracket.  From
$e=\mu/(L+n+k)$,
\[
  1+e=\frac{(L+n+k)+(L+n-k)}{L+n+k}=\frac{2(L+n)}{L+n+k},
  \qquad
  \frac{\mu}{2}-n=\frac{L+n-k}{2}-n=\frac{L-n-k}{2},
\]
so
\[
  \Bigl(\frac{\mu}{2}-n\Bigr)(1+e)
  =\frac{L-n-k}{2}\cdot\frac{2(L+n)}{L+n+k}
  =\frac{(L-n-k)(L+n)}{L+n+k},
\]
and therefore
\[
  k+\Bigl(\frac{\mu}{2}-n\Bigr)(1+e)
  =\frac{k(L+n+k)+(L-n-k)(L+n)}{L+n+k}
  =\frac{L^2+k^2-n^2}{L+n+k},
\]
the numerator because
$k(L+n+k)+(L-n-k)(L+n)=kL+kn+k^2+L^2+Ln-nL-n^2-kL-kn=L^2+k^2-n^2$.
By (I), $L^2=-pnL-(n^2-k^2)(p-1)$, so
\[
  L^2+k^2-n^2=-pnL-(n^2-k^2)(p-1)+k^2-n^2
  =-pnL+(k^2-n^2)\bigl[(p-1)+1\bigr]=p\bigl(k^2-n^2-nL\bigr).
\]
Since $p>1>0$ we have $c=\tfrac12-\tfrac1p=\dfrac{p-2}{2p}$, and hence
\begin{equation}\tag{7}
  c\Bigl[k+\Bigl(\frac{\mu}{2}-n\Bigr)(1+e)\Bigr]
  =\frac{p-2}{2p}\cdot\frac{p\bigl(k^2-n^2-nL\bigr)}{L+n+k}
  =\frac{(p-2)\bigl(k^2-n^2-nL\bigr)}{2(L+n+k)} .
\end{equation}
So $(\alpha)$, which reads $\mu/2=(L+n-k)/2$ on the left, is equivalent --- after
multiplying by $2(L+n+k)>0$ --- to
\begin{equation}\tag{8}
  (p-2)\bigl(k^2-n^2-nL\bigr)=(L+n-k)(L+n+k)=(L+n)^2-k^2 .
\end{equation}
We verify (8).  Using (I) once more,
\begin{align*}
  (L+n)^2-k^2&=L^2+2nL+n^2-k^2\\
  &=-pnL-(n^2-k^2)(p-1)+2nL+n^2-k^2\\
  &=nL(2-p)+(k^2-n^2)(p-1)+(n^2-k^2)\\
  &=nL(2-p)+(k^2-n^2)\bigl[(p-1)-1\bigr]
  =nL(2-p)+(k^2-n^2)(p-2),
\end{align*}
whereas the left-hand side of (8) is
\[
  (p-2)(k^2-n^2)-(p-2)nL=(k^2-n^2)(p-2)+nL(2-p).
\]
The two agree, so (8) holds, and with it $(\alpha)$ and, by Step 6, $(\beta)$.

\emph{Step 8: the pointwise form of (18).}  By Steps 5--7, equation (18) holds
at every $\xi\neq0$; that is,
\begin{equation}\tag{9}
  \frac{\partial T_k}{\partial\bar\xi}(\xi)
  =c\Bigl[\frac{\xi}{\bar\xi}\frac{\partial T_k}{\partial\xi}(\xi)
    +\frac{\bar\xi^{\,n}}{\xi^{n}}
      \overline{\frac{\partial T_k}{\partial\xi}(\xi)}\Bigr]
  \qquad(\xi\neq0),
\end{equation}
the derivatives being the pointwise operators of \noderef{DZ} and
\noderef{DZbar}.  In particular (9) holds at every $\xi\in\mathbb{B}$ with
$\xi\neq0$.

\emph{Step 9: growth bounds for $T_k$ and its first derivatives.}  Put
$M=(1+|e|)|A_k|$ and let $\xi\neq0$, $r=|\xi|$.  Since $|\bar\xi|=|\xi|=r$ and
$e$ is real,
\[
  |F(\xi)|\le|A_k|r^{k}+|e|\,|A_k|r^{k}=M\,r^{k},
\]
so by (0), $\bigl||\xi|^{\mu}\bigr|=r^{\mu}$, $|\xi^{-n}|=r^{-n}$ and
$k+\mu-n=L$ (Step 1),
\begin{equation}\tag{10}
  |T_k(\xi)|\le M\,r^{k+\mu-n}=M\,r^{L}.
\end{equation}
By the two displays of Step 4, with $|\xi^{k-1}|=|\bar\xi^{\,k-1}|=r^{k-1}$,
$|1/\xi|=|1/\bar\xi|=r^{-1}$ and the bound on $|F|$,
\[
  \Bigl|\frac{\partial T_k}{\partial\xi}(\xi)\Bigr|
  \le r^{\mu-n}\Bigl[k|A_k|r^{k-1}
    +\Bigl|\frac{\mu}{2}-n\Bigr|M\,r^{k-1}\Bigr],
  \qquad
  \Bigl|\frac{\partial T_k}{\partial\bar\xi}(\xi)\Bigr|
  \le r^{\mu-n}\Bigl[k|e||A_k|r^{k-1}
    +\frac{|\mu|}{2}M\,r^{k-1}\Bigr],
\]
so, with $M_1=M\bigl(k+\bigl|\tfrac{\mu}{2}-n\bigr|+\tfrac{|\mu|}{2}\bigr)$ and
$\mu-n+k-1=L-1$,
\begin{equation}\tag{11}
  \Bigl|\frac{\partial T_k}{\partial\xi}(\xi)\Bigr|\le M_1r^{L-1},
  \qquad
  \Bigl|\frac{\partial T_k}{\partial\bar\xi}(\xi)\Bigr|\le M_1r^{L-1}
  \qquad(\xi\neq0),
\end{equation}
using $|A_k|\le M$ and $|e|\,|A_k|\le M$.

Write $\partial_0=D^{[0]}$ and $\partial_1=D^{[1]}$ for the classical partial
derivatives of \noderef{PartialWord}, so that
$\partial_0f(\xi)=\mathrm{d}f(\xi)[1]$ and $\partial_1f(\xi)=\mathrm{d}f(\xi)[i]$.
By \noderef{DZ} and \noderef{DZbar},
$\partial T_k/\partial\xi=\tfrac12(\partial_0T_k-i\,\partial_1T_k)$ and
$\partial T_k/\partial\bar\xi=\tfrac12(\partial_0T_k+i\,\partial_1T_k)$ at
every point, whence
\begin{equation}\tag{12}
  \partial_0T_k=\frac{\partial T_k}{\partial\xi}
    +\frac{\partial T_k}{\partial\bar\xi},
  \qquad
  \partial_1T_k=i\Bigl(\frac{\partial T_k}{\partial\xi}
    -\frac{\partial T_k}{\partial\bar\xi}\Bigr)
\end{equation}
at every point.  With (11), and with $L\ge0$ from (1) so that
$r^{L-1}=r^{L}r^{-1}\le r^{-1}$ for $0<r\le1$,
\begin{equation}\tag{13}
  |T_k(\xi)|\le M,\qquad |\partial_jT_k(\xi)|\le2M_1\,|\xi|^{-1}
  \qquad(0<|\xi|<1,\ j=0,1),
\end{equation}
the first by (10) and $r^{L}\le1$.

\emph{Step 10: the weak first partial derivatives of $T_k$ on $\mathbb{B}$.}
Fix $j\in\{0,1\}$.  We verify the hypotheses of
\noderef{WeakPartialAcrossPoint} with $\Omega=\mathbb{B}$, $S=\{0\}$, this
$j$, $g=T_k$, $v=\partial_jT_k$ and $\beta=0$.

The set $\mathbb{B}$ is open and $S=\{0\}\subseteq\mathbb{B}$; $S\cap\kappa$ is
finite for every compact $\kappa$, $S$ being a singleton; and
$0\le\beta=0<1$.

$g$ is of class $C^1$ on $\mathbb{B}-S$: the set $\mathbb{B}-\{0\}$ is open and
$T_k$ is $C^\infty$ at each of its points by (a), and a function of class
$C^\infty$ on an open set is of class $C^1$ there.  And $D^{[j]}g=v$ on
$\mathbb{B}-S$ holds by the definition of $v$.

$g\in L^1_{\mathrm{loc}}(\mathbb{B})$ in the sense of \noderef{MemLpLocOn}:
$T_k$ is continuous on the open set $\mathbb{B}-\{0\}$ by (a), hence
measurable there, and $\{0\}$ is a null set, so $T_k$ is measurable on
$\mathbb{B}$ for the completed Lebesgue measure; and by (13), for every
compact $\kappa\subseteq\mathbb{B}$,
$\int_\kappa|T_k|\,d\sigma\le M\,\sigma(\kappa)<\infty$, $\sigma(\kappa)$
being finite because $\kappa$ is bounded.

$v\in L^1_{\mathrm{loc}}(\mathbb{B})$: measurability is the same argument,
$\partial_jT_k$ being continuous on $\mathbb{B}-\{0\}$ by (a).  For
integrability, apply \noderef{PowerIntegrableNearZero} with $\alpha=1$,
$\delta=1$ and $q=1$; its hypotheses $0<\alpha$, $\alpha\le1$, $0<\delta$,
$1\le q$ and $q<2/(2-\alpha)=2$ all hold, and its conclusion is that
$\xi\mapsto|\xi|^{(\alpha-2)q}=|\xi|^{-1}$ is integrable on
$\mathbb{B}=\{|\xi|<1\}$.  By (13), $|v|\le2M_1|\xi|^{-1}$ on
$\mathbb{B}-\{0\}$, hence almost everywhere on $\mathbb{B}$, so
$\int_{\mathbb{B}}|v|\,d\sigma<\infty$ and a fortiori
$\int_\kappa|v|\,d\sigma<\infty$ for every compact $\kappa\subseteq\mathbb{B}$.

The growth hypothesis: at the only point $0$ of $S$, take the constants $M$
and $\rho=1$; by (13), $\|T_k(\xi)\|\le M=M\,|\xi-0|^{-0}$ whenever
$0<|\xi-0|<1$.

Therefore $\partial_jT_k$ is the distributional $\partial_j$-derivative of
$T_k$ on $\mathbb{B}$, in the sense of \noderef{IsWeakPartialOn}, for
$j=0,1$.

\emph{Step 11: conclusion.}  Put
$a=\partial T_k/\partial\xi$ and $b=\partial T_k/\partial\bar\xi$, and
$f_x=\partial_0T_k$, $f_y=\partial_1T_k$.  The three clauses of
\noderef{IsWeakWirtingerOn} hold for $T_k$, $a$, $b$ on $\mathbb{B}$ with this
$f_x,f_y$: clause (a) is the local integrability of $T_k$, $f_x$ and $f_y$
established in Step 10; clause (b) is the conclusion of Step 10; and clause
(c) is the identity
$a=\tfrac12(f_x-if_y)$, $b=\tfrac12(f_x+if_y)$, valid at every point of
$\mathbb{C}$ by \noderef{DZ} and \noderef{DZbar} as recorded before (12).  So
$a$ and $b$ are weak Wirtinger derivatives of $T_k$ on $\mathbb{B}$.

By Step 8 the identity (9), which is exactly
\[
  b(\xi)=c\Bigl[\frac{\xi}{\bar\xi}a(\xi)
    +\frac{\bar\xi^{\,n}}{\xi^{n}}\overline{a(\xi)}\Bigr],
\]
holds at every $\xi\in\mathbb{B}$ with $\xi\neq0$, hence at almost every such
$\xi$.  With the pair $a,b$ this is precisely the assertion that $T_k$ solves
(18) on $\mathbb{B}$ in the sense of \noderef{SolvesEq18}.  This is (c), and
the proof is complete.
\end{proof}
